% from ega4-2-fr.tex
\oldpage[IV-2]{110}

On notera que l'ensemble $Z$ n'est pas nécessairement constructible: on a un
exemple de ce fait en prenant $X=\operatorname{Spec}(A)$, où $A$ est un
anneau intègre noethérien ayant une infinité d'idéaux maximaux, et $Z$ le
complémentaire du point générique de $X$.

Si $Z$ est fermé et si $i:X-Z\to X$ est l'injection canonique, on a
\begin{equation}
\mathcal H^0_{X/Z}(\mathcal F)=i_*(\mathcal F|X-Z)
\tag{5.9.1.3}\label{IV.2.s5.9.1.3-fr}
\end{equation}
et en particulier, pour $Z=\varnothing$,
$\mathcal H^0_{X/Z}(\mathcal F)=\mathcal F$.

\begin{proposition}[5.9.2]
\label{IV.2.s5.9.2-fr}
\begin{enumerate}
\item[(i)] Le foncteur
$\mathcal F\mapsto\mathcal H^0_{X/Z}(\mathcal F)$ est exact à gauche.
\item[(ii)] Si $\mathcal F$ est quasi-cohérent, il en est de même de
$\mathcal H^0_{X/Z}(\mathcal F)$.
\end{enumerate}
\end{proposition}

L'assertion (i) résulte de la définition (5.9.1.1), du fait que
$(i_\alpha)_*$ est un foncteur exact à gauche et de ce que la limite
inductive préserve l'exactitude dans la catégorie des $\mathcal O_X$-Modules.
L'assertion (ii) résulte de (I, 9.2.2) et du fait qu'une limite inductive de
$\mathcal O_X$-Modules quasi-cohérents est quasi-cohérente (I, 1.3.9).

\begin{remark}[5.9.3]
\label{IV.2.s5.9.3-fr}
Si $\mathcal F$ est une $\mathcal O_X$-Algèbre, il en est de même de
$\mathcal H^0_{X/Z}(\mathcal F)$ $(0_{\mathrm I},4.2.4)$; en particulier
$\mathcal H^0_{X/Z}(\mathcal O_X)$ est une $\mathcal O_X$-Algèbre
quasi-cohérente, et pour tout $\mathcal O_X$-Module $\mathcal F$,
$\mathcal H^0_{X/Z}(\mathcal F)$ est un
$\mathcal H^0_{X/Z}(\mathcal O_X)$-Module, qui est quasi-cohérent si
$\mathcal F$ est quasi-cohérent (I, 9.6.1). Plus particulièrement,
supposons que $X=\operatorname{Spec}(A)$, où $A$ est intègre et noethérien;
alors $\mathcal H^0_{X/Z}(\mathcal O_X)$ est la $\mathcal O_X$-Algèbre
$\widetilde B$, où
\begin{equation}
B=\bigcap_{\mathfrak p\in X-Z}A_{\mathfrak p}.
\tag{5.9.3.1}\label{IV.2.s5.9.3.1-fr}
\end{equation}
Cela résulte en effet de (5.9.1.2) et de (I, 8.2.1.1).
\end{remark}

\begin{proposition}[5.9.4]
\label{IV.2.s5.9.4-fr}
Soient $X$ un préschéma localement noethérien, $Z$ une partie de $X$ stable
par spécialisation, $X'$ un préschéma localement noethérien,
$f:X'\to X$ un morphisme plat. Alors $Z'=f^{-1}(Z)$ est stable par
spécialisation et pour tout $\mathcal O_X$-Module quasi-cohérent $\mathcal F$,
on a un isomorphisme canonique
\begin{equation}
f^*(\mathcal H^0_{X/Z}(\mathcal F))
\xrightarrow{\sim}\mathcal H^0_{X'/Z'}(f^*(\mathcal F)).
\tag{5.9.4.1}\label{IV.2.s5.9.4.1-fr}
\end{equation}
\end{proposition}

En effet, avec les notations de (5.9.1), $Z'_\alpha=f^{-1}(Z_\alpha)$ est
fermé dans $X'$ et $Z'$ est réunion des $Z'_\alpha$; en outre,
$(i_\alpha)_{(X')}$ est l'injection canonique
$i'_\alpha:U'_\alpha\to X'$, si
$U'_\alpha=X'-Z'_\alpha=f^{-1}(U_\alpha)$. Comme $f$ est plat, on sait
(2.3.1) que l'homomorphisme canonique
\[
f^*((i_\alpha)_*(\mathcal F|U_\alpha))
\longrightarrow(i'_\alpha)_*(f^*(\mathcal F)|U'_\alpha)
\]
est bijectif; comme, pour $\alpha\leq\beta$, le diagramme
\[
\xymatrix{
f^*((i_\alpha)_*(\mathcal F|U_\alpha))
  \ar[r]^{\sim}\ar[d] &
(i'_\alpha)_*(f^*(\mathcal F)|U'_\alpha)\ar[d]\\
f^*((i_\beta)_*(\mathcal F|U_\beta))
  \ar[r]^{\sim} &
(i'_\beta)_*(f^*(\mathcal F)|U'_\beta)
}
\]
est commutatif, on a donc, en passant à la limite, un isomorphisme canonique
\[
\varinjlim_\alpha
f^*((i_\alpha)_*(\mathcal F|U_\alpha))
\xrightarrow{\sim}\mathcal H^0_{X'/Z'}(f^*(\mathcal F)).
\]

